% !TeX root=main.tex
% در این فایل، عنوان پایان‌نامه، مشخصات خود، متن تقدیمی‌، ستایش، سپاس‌گزاری و چکیده پایان‌نامه را به فارسی، وارد کنید.
% توجه داشته باشید که جدول حاوی مشخصات پروژه/پایان‌نامه/رساله و همچنین، مشخصات داخل آن، به طور خودکار، درج می‌شود.
%%%%%%%%%%%%%%%%%%%%%%%%%%%%%%%%%%%%
% دانشگاه خود را وارد کنید
\university{علم و صنعت ایران}
% دانشکده، آموزشکده و یا پژوهشکده  خود را وارد کنید
\faculty{دانشکده مهندسی کامپیوتر}
% گروه آموزشی خود را وارد کنید
\department{گروه شبکه}
% گروه آموزشی خود را وارد کنید
\subject{مهندسی کامپیوتر}
% گرایش خود را وارد کنید
\field{شبکه}
% عنوان پایان‌نامه را وارد کنید
\title{همكاري در توسعه يك برنامه كاربردي موبايل براي كمك به رشد گفتار كودكان زير 3 سال}
% نام استاد(ان) راهنما را وارد کنید
\firstsupervisor{خانم دکتر زینب موحدی}
% \secondsupervisor{استاد راهنمای دوم}
% نام استاد(دان) مشاور را وارد کنید. چنانچه استاد مشاور ندارید، دستور پایین را غیرفعال کنید.
% \firstadvisor{استاد مشاور اول}
%\secondadvisor{استاد مشاور دوم}
% نام دانشجو را وارد کنید
\name{محمدرضا}
% نام خانوادگی دانشجو را وارد کنید
\surname{صائمي پور}
% شماره دانشجویی دانشجو را وارد کنید
\studentID{۹۹۵۲۱۳۸۸}
% تاریخ پایان‌نامه را وارد کنید
\thesisdate{خرداد ۱۴۰۵}
% به صورت پیش‌فرض برای پایان‌نامه‌های کارشناسی تا دکترا به ترتیب از عبارات «پروژه»، «پایان‌نامه» و »رساله» استفاده می‌شود؛ اگر  نمی‌پسندید هر عنوانی را که مایلید در دستور زیر قرار داده و آنرا از حالت توضیح خارج کنید.
%\projectLabel{پایان‌نامه}

% به صورت پیش‌فرض برای عناوین مقاطع تحصیلی کارشناسی تا دکترا به ترتیب از عبارات «کارشناسی»، «کارشناسی ارشد» و »دکترا» استفاده می‌شود؛ اگر  نمی‌پسندید هر عنوانی را که مایلید در دستور زیر قرار داده و آنرا از حالت توضیح خارج کنید.
%\degree{}

\firstPage
\besmPage
\davaranPage

%\vspace{.5cm}
% در این قسمت اسامی اساتید راهنما، مشاور و داور باید به صورت دستی وارد شوند
\renewcommand{\arraystretch}{1.6}
\begin{center}
\small % or \footnotesize for even smaller font
\begin{tabular}{| p{12mm} | p{18mm} | p{.25\textwidth} |p{14mm}|p{.2\textwidth}|c|}
\hline
ردیف	& سمت & نام و نام خانوادگی & مرتبه \newline دانشگاهی &	دانشگاه یا مؤسسه &	امضـــــــــــــا\\
\hline
۱  &	استاد راهنما				 & خانم دکتر زینب موحدی & استادیار & علم و صنعت ایران &  \\
\hline
۲ &	استاد داور			 & دکتر حسن نادری & دانشیار & علم و صنعت ایران & \\
\hline
\end{tabular}
\normalsize % Reset to normal size after the table
\end{center}


\esalatPage
\mojavezPage


% چنانچه مایل به چاپ صفحات «تقدیم»، «نیایش» و «سپاس‌گزاری» در خروجی نیستید، خط‌های زیر را با گذاشتن ٪  در ابتدای آنها غیرفعال کنید.
 % پایان‌نامه خود را تقدیم کنید!
\newpage
\thispagestyle{empty}

\begin{flushright}
	
	{\Large \textbf{تقدیم به}}\\[10mm]
	
	\setstretch{1.4}
	{\large
		مادر و همسر عزیزم،\\
		که با عشق بی‌کران، صبوری و حمایت‌های همیشگی خود، همواره پشتیبان و الهام‌بخش من در مسیر زندگی و تحصیل بوده‌اند.\\[6mm]
		
		استادان گرامی‌ام،\\
		که با آموزش‌های ارزشمند، نگاه علمی دقیق و راهنمایی‌های دلسوزانه‌ی خود، زمینه‌ساز رشد فکری و پیشرفت تحصیلی‌ام شدند.
	}
	
\end{flushright}

\vfill


% سپاس‌گزاری
\begin{acknowledgementpage}
با نهایت قدردانی و احترام، از استاد ارجمندم، خانم دکتر زینب موحدی، صمیمانه سپاسگزارم که با دانش عمیق، دیدگاه‌های راهگشا و راهنمایی‌های دلسوزانه‌ی خود در تمامی مراحل انجام این پژوهش همراه و پشتیبان من بودند. بی‌تردید، این اثر حاصل حمایت‌های علمی، اعتماد و هدایت‌های ارزشمند ایشان است و بدون حضور مؤثر و حمایت‌های بی‌دریغشان، دستیابی به نتایج این تحقیق میسر نبود.
% با استفاده از دستور زیر، امضای شما، به طور خودکار، درج می‌شود.
\signature 
\end{acknowledgementpage}
%%%%%%%%%%%%%%%%%%%%%%%%%%%%%%%%%%%%
% کلمات کلیدی پایان‌نامه را وارد کنید
\keywords{رشد گفتار کودک، یادگیری واژگان، آموزش دیداری‌ـ‌شنیداری، بازی‌وارسازی، نرم‌افزار همراه، کاتلین چندسکویی}

%چکیده پایان‌نامه را وارد کنید، برای ایجاد پاراگراف جدید از \\ استفاده کنید. اگر خط خالی دشته باشید، خطا خواهید گرفت.
\fa-abstract{سه سال نخست زندگی، دوره‌ای تعیین‌کننده در شکل‌گیری مهارت‌های گفتاری و گسترش دامنه واژگان کودک است و فراهم‌کردن تعامل زبانی مستمر و متناسب با توانایی او می‌تواند به این فرایند کمک کند. با این حال، والدین همواره به محتوایی ساخت‌یافته، جذاب و قابل‌شخصی‌سازی برای تمرین روزانه با کودک دسترسی ندارند. هدف این پروژه، همکاری در طراحی و توسعه «نوواژه» به‌عنوان یک برنامه کاربردی همراه برای پشتیبانی از یادگیری واژگان و ایجاد فرصت‌های تعاملی میان والد و کودک زیر سه سال است. در این برنامه، واژه‌های فارسی در دسته‌بندی‌های موضوعی و با استفاده از تصویر و صوت ارائه می‌شوند. کودک پس از شنیدن تلفظ هر واژه، تصویر متناظر را انتخاب می‌کند و بازخورد دیداری و شنیداری دریافت می‌کند. پرسش‌ها با اولویت‌دادن به واژه‌هایی که هنوز آموخته نشده‌اند تولید می‌شوند و مجموعه‌ای از بازی‌های تعاملی نیز با هدف حفظ انگیزه و تقویت توجه، حافظه و تشخیص دیداری در اختیار کودک قرار می‌گیرد. بخش ویژه والدین امکان ایجاد و مدیریت نمایه فرزندان، افزودن دسته‌بندی و واژه دلخواه با تصویر و صدای ضبط‌شده، تعیین زمان مجاز استفاده و فعال‌سازی حالت کودک را فراهم می‌کند. همچنین، پاسخ‌های کودک ثبت می‌شوند تا شاخص‌هایی مانند تعداد تلاش‌ها، پاسخ‌های درست، دقت کلی و پیشرفت در هر دسته‌بندی قابل محاسبه و پیگیری باشد. نوواژه با معماری کارخواه‌ـ‌کارساز توسعه یافته است؛ رابط کاربری مشترک اندروید و آی‌اواس با کاتلین و کامپوز چندسکویی پیاده‌سازی شده و کارساز مبتنی بر کتور، مدیریت محتوا، احراز هویت و نگهداری روند یادگیری را بر عهده دارد. حاصل این پروژه، یک نمونه کاربردی چندسکویی و توسعه‌پذیر است که آموزش دیداری‌ـ‌شنیداری، بازی‌وارسازی، شخصی‌سازی محتوا و پایش پیشرفت را در محیطی کودک‌محور یکپارچه می‌کند و می‌تواند به‌عنوان ابزار کمکی در کنار تعامل مستقیم والد با کودک به کار رود.}


\abstractPage

\newpage\clearpage